% Part: computability
% Chapter: computability-theory
% Section: computable-sets

\documentclass[../../../include/open-logic-section]{subfiles}

\begin{document}

\olfileid{cmp}{thy}{cps}
\olsection{المجموعات القابلة للحساب}

يمكننا توسيع مفهوم قابلية الحساب من الدوال القابلة للحساب إلى
المجموعات القابلة للحساب:

\begin{defn}
  لتكن $S$ مجموعة من الأعداد الطبيعية. تكون $S$ \emph{قابلة للحساب}
  إذا وفقط إذا كانت دالتها المميزة~$\Char{S}$ قابلة للحساب. وبعبارة
  أخرى، تكون $S$ قابلة للحساب إذا وفقط إذا كانت الدالة
\[
\Char{S}(x) =
\begin{cases}
1 & \text{إذا كان $x \in S$} \\
0 & \text{في غير ذلك}
\end{cases}
\]
قابلة للحساب. وبالمثل، تكون العلاقة
$R(x_0,\dots,x_{k-1})$ قابلة للحساب إذا وفقط إذا كانت دالتها المميزة
كذلك.

وتسمى المجموعات والعلائق القابلة للحساب أيضًا \emph{قابلة للقرار}.
\end{defn}

\begin{explain}
لاحظ أن لدينا الآن عدة مفاهيم لقابلية الحساب: للدوال الجزئية، وللدوال،
وللمجموعات. فلا تخلط بينها!{} \iftag{TMs}{تعيد آلة تورنغ التي تحسب
دالة جزئية خرج الدالة عند قيم الدخل التي تكون الدالة معرفة عندها؛ أما
آلة تورنغ التي تحسب مجموعة فتعيد إما $1$ وإما~$0$، بعد أن تقرر ما إذا
كانت قيمة الدخل في المجموعة.}{}
\end{explain}

\end{document}
